\subsection{Expectation and Integration in Probability}
% Expectation, moments, variance, covariance, Jensen, Markov/Chebyshev inequalities, integration as probabilistic averaging.
% Source integration: attached notation pages define expectation, variance, covariance, and expectation of functions through the distribution function.

Expectation is the averaging operation attached to a probability law. In the notation of Section \ref{def:probability-space-random-variable-law}, if $X$ has law $F_X$, then integration against $F_X$ is the common notation that covers both sums for discrete variables and ordinary integrals for variables with densities.

\begin{definition} \label{def:expectation-discrete-continuous-and-functionals} \hlt{Expectation Through the Law}\\
Let $X$ be real-valued with distribution function $F_X$.
\begin{enumerate}[label=(\roman*)]
\setlength{\itemsep}{0pt}
\item \hlt{Expectation}: if $\int_{\R}|x|\,dF_X(x)<\infty$, then $\mathbb{E}(X)\defeq\int_{\R}x\,dF_X(x)$.
\item \hlt{Discrete form}: if $X$ has support $S$, mass function $f_X$, $\mathbb{E}(X)=\sum\limits_{x\in S}xf_X(x)$ when $\sum\limits_{x\in S}|x|f_X(x)<\infty$.
\item \hlt{Density form}: if $X$ has density $f_X$, then $\mathbb{E}(X)=\int_{\R}xf_X(x)\,dx$ when $\int_{\R}|x|f_X(x)\,dx<\infty$.
\item \hlt{Function form}: for measurable $r:\R\to\R$, if $\int_{\R}|r(x)|\,dF_X(x)<\infty$, then $\mathbb{E}(r(X))\defeq\int_{\R}r(x)\,dF_X(x)$.
\end{enumerate}
\end{definition}

\begin{example} \label{ex:expectation-from-two-coin-count} \hlt{Expectation of the Two-Toss Head Count}\\
Let $\Omega=\{H,T\}^2$ with $\mathbb{P}(\{\omega\})=\frac{1}{4}$, and let $X(\omega)$ be the number of heads in $\omega$.\\
Then $f_X(k)=\binom{2}{k}2^{-2}$ for $k=0,1,2$, so $\mathbb{E}(X)=\sum\limits_{k=0}^{2}k f_X(k)=0\cdot\frac{1}{4}+1\cdot\frac{1}{2}+2\cdot\frac{1}{4}=1$.
\end{example}

\begin{definition} \label{def:variance-covariance-and-standard-deviation} \hlt{Variance and Covariance}
\begin{enumerate}[label=(\roman*)]
\setlength{\itemsep}{0pt}
\item The \hlt{variance} of $X$ is $\operatorname{Var}(X)\defeq\mathbb{E}[(X-\mathbb{E}X)^2]$, assuming $\mathbb{E}(X^2)<\infty$.
\item The \hlt{standard deviation} of $X$ is $\operatorname{sd}(X)\defeq\sqrt{\operatorname{Var}(X)}$.
\item The \hlt{covariance} of $X,Y$ is $\operatorname{Cov}(X,Y)\defeq\mathbb{E}[(X-\mathbb{E}X)(Y-\mathbb{E}Y)]$, assuming $\mathbb{E}(X^2),\mathbb{E}(Y^2)<\infty$.
\end{enumerate}
\end{definition}

\begin{proposition} \label{prop:indicator-expectation-probability} \hlt{Indicators Turn Probabilities into Expectations}\\
For $A\in\mathcal{F}$, let $\mathbf{1}_A:\Omega\to\{0,1\}$ be $\mathbf{1}_A(\omega)=1$ for $\omega\in A$ and $\mathbf{1}_A(\omega)=0$ for $\omega\notin A$. Then $\mathbb{E}(\mathbf{1}_A)=\mathbb{P}(A)$.
\begin{proof}
Since $\Omega=A\sqcup A^c$ and $\mathbf{1}_A=1$ on $A$, $0$ on $A^c$, $\mathbb{E}(\mathbf{1}_A)=1\cdot\mathbb{P}(A)+0\cdot\mathbb{P}(A^c)=\mathbb{P}(A)$.\\
\hlt{Structural Commentary}: Indicators encode events as $0$-$1$ random variables.\\
\hlt{Key Insight}: Probabilities are expectations of indicators.
\end{proof}
\end{proposition}

\begin{remark} \label{rem:expectation-section-source-scope} \hlt{Law-Based Integration}\\
For measurable $r$, $\mathbb{E}(r(X))$ depends only on the law of $X$: $\mathbb{E}(r(X))=\int_{\R}r(x)\,dF_X(x)$ whenever integrable.\\
Discrete laws give $\sum\limits_x r(x)f_X(x)$; density laws give $\int_{\R}r(x)f_X(x)\,dx$.
\end{remark}

\begin{example} \label{ex:first-expectation-computations} \hlt{First Expectation Computations}
\begin{enumerate}[label=(\roman*)]
\setlength{\itemsep}{0pt}
\item If $X\sim\operatorname{Bernoulli}(p)$, then $\mathbb{E}X=0(1-p)+1p=p$.
\item If $X\sim\operatorname{Uniform}(-1,3)$, then $\mathbb{E}X=(\frac{1}{4})\int_{-1}^{3}x\,dx=1$.
\item If $X\sim\operatorname{Cauchy}(0,1)$, then $\mathbb{E}|X|=(\frac{2}{\pi})\int_0^\infty x(1+x^2)^{-1}\,dx=\infty$, so $\mathbb{E}X$ does not exist.
\end{enumerate}
\end{example}

\begin{theorem} \label{thm:law-of-unconscious-statistician} \hlt{Expectation of a Function}\\
Let $r$ be measurable and assume the displayed expectations exist.
\begin{enumerate}[label=(\roman*)]
\setlength{\itemsep}{0pt}
\item \hlt{One variable}: for real-valued $X$, $\mathbb{E}(r(X))=\int_{\R}r(x)\,dF_X(x)$.
\item \hlt{Several variables}: for $(X,Y)$ with joint law $F_{X,Y}$, $\mathbb{E}(r(X,Y))=\int_{\R^2}r(x,y)\,dF_{X,Y}(x,y)$.
\end{enumerate}
\begin{proof}
\begin{enumerate}[label=(\roman*)]
\setlength{\itemsep}{0pt}
\item Since $\mathcal{L}(X)=X_\ast\mathbb{P}$, $\mathbb{E}(r(X))=\int_\Omega r(X(\omega))\,d\mathbb{P}(\omega)=\int_{\R}r(x)\,d\mathcal{L}(X)(x)=\int_{\R}r(x)\,dF_X(x)$.
\item Let $W=(X,Y)$. Since $\mathcal{L}(W)=W_\ast\mathbb{P}=F_{X,Y}$,\\
$\mathbb{E}(r(X,Y))=\int_\Omega r(W(\omega))\,d\mathbb{P}(\omega)=\int_{\R^2}r(x,y)\,d\mathcal{L}(W)(x,y)=\int_{\R^2}r(x,y)\,dF_{X,Y}(x,y)$.
\end{enumerate}
\hlt{Structural Commentary}: Expectations of transformations integrate over the transformed law.\\
\hlt{Key Insight}: Pushforward laws let functions move from random variables to their values.
\end{proof}
\end{theorem}

\begin{example} \label{ex:lotus-transform-examples} \hlt{Using the Function Rule}
\begin{enumerate}[label=(\roman*)]
\setlength{\itemsep}{0pt}
\item If $X\sim\operatorname{Uniform}(0,1)$ and $Y=e^X$, then $\mathbb{E}Y=\mathbb{E}(e^X)=\int_0^1e^x\,dx=e-1$.
\item If $X\sim\operatorname{Uniform}(0,1)$ and $Y=\max\{X,1-X\}$, then $\mathbb{E}Y=\int_0^{\frac{1}{2}}(1-x)\,dx+\int_{\frac{1}{2}}^{1}x\,dx=\frac{3}{4}$.
\item If $(X,Y)\sim\operatorname{Uniform}([0,1]^2)$ and $Z=X^2+Y^2$, then $\mathbb{E}Z=\int_0^1\int_0^1(x^2+y^2)\,dx\,dy=\frac{2}{3}$.
\end{enumerate}
\end{example}

\begin{definition} \label{def:moments-correlation-and-sample-statistics} \hlt{Moments, Correlation, and Sample Statistics}
\begin{enumerate}[label=(\roman*)]
\setlength{\itemsep}{0pt}
\item The \hlt{$k$th moment} of $X$ is $\mathbb{E}(X^k)$, assuming $\mathbb{E}|X|^k<\infty$.
\item The \hlt{correlation} of $X,Y$ is $\operatorname{Corr}(X,Y)\defeq\frac{\operatorname{Cov}(X,Y)}{\sqrt{\operatorname{Var}(X)}\sqrt{\operatorname{Var}(Y)}}$, assuming $\operatorname{Var}(X),\operatorname{Var}(Y)>0$.
\item The \hlt{sample mean} of $X_1,\ldots,X_n$ is $\overline{X}_n\defeq n^{-1}\sum\limits_{i=1}^{n}X_i$.
\item The \hlt{sample variance} of $X_1,\ldots,X_n$ is $S_n^2\defeq(n-1)^{-1}\sum\limits_{i=1}^{n}(X_i-\overline{X}_n)^2$.
\end{enumerate}
\end{definition}

\begin{theorem} \label{thm:expectation-variance-covariance-rules} \hlt{Expectation, Variance, and Covariance Algebra}\\
Assume all displayed quantities exist and $a_i,a,b\in\R$.
\begin{enumerate}[label=(\roman*)]
\setlength{\itemsep}{0pt}
\item \hlt{Linearity}: $\mathbb{E}(\sum\limits_i a_iX_i)=\sum\limits_i a_i\mathbb{E}X_i$.
\item \hlt{Independent product}: if $X_1,\ldots,X_n$ are independent, then $\mathbb{E}(\prod\limits_i X_i)=\prod\limits_i\mathbb{E}X_i$.
\item \hlt{Variance identities}: $\operatorname{Var}(X)=\mathbb{E}(X^2)-(\mathbb{E}X)^2$ and $\operatorname{Var}(aX+b)=a^2\operatorname{Var}(X)$.
\item \hlt{Covariance identity}: $\operatorname{Cov}(X,Y)=\mathbb{E}(XY)-\mathbb{E}X\,\mathbb{E}Y$.
\item \hlt{Variance of a linear combination}: $\operatorname{Var}(\sum\limits_i a_iX_i)=\sum\limits_i a_i^2\operatorname{Var}(X_i)+2\sum\limits_{i<j}a_ia_j\operatorname{Cov}(X_i,X_j)$.
\end{enumerate}
\begin{proof}
\begin{enumerate}[label=(\roman*)]
\setlength{\itemsep}{0pt}
\item $\mathbb{E}(\sum\limits_i a_iX_i)=\int_\Omega\sum\limits_i a_iX_i\,d\mathbb{P}=\sum\limits_i a_i\int_\Omega X_i\,d\mathbb{P}=\sum\limits_i a_i\mathbb{E}X_i$.
\item By independence, $\mathcal{L}(X_1,\ldots,X_n)=\prod\limits_i F_{X_i}$; $\mathbb{E}(\prod\limits_i X_i)=\int_{\R^n}\prod\limits_i x_i\,d\prod\limits_i F_{X_i}(x_i)=\prod\limits_i\int_\R x_i\,dF_{X_i}(x_i)=\prod\limits_i\mathbb{E}X_i$.
\item $\operatorname{Var}(X)=\mathbb{E}(X^2-2X\mathbb{E}X+(\mathbb{E}X)^2)=\mathbb{E}(X^2)-(\mathbb{E}X)^2$.\\
Also $\operatorname{Var}(aX+b)=\mathbb{E}[(aX+b-a\mathbb{E}X-b)^2]=a^2\operatorname{Var}(X)$.
\item $\operatorname{Cov}(X,Y)=\mathbb{E}(XY-X\mathbb{E}Y-Y\mathbb{E}X+\mathbb{E}X\mathbb{E}Y)=\mathbb{E}(XY)-\mathbb{E}X\,\mathbb{E}Y$.
\item With $S=\sum\limits_i a_iX_i$, (iv) gives\\
$\operatorname{Var}(S)=\operatorname{Cov}(S,S)=\sum\limits_{i,j}a_ia_j\operatorname{Cov}(X_i,X_j)=\sum\limits_i a_i^2\operatorname{Var}(X_i)+2\sum\limits_{i<j}a_ia_j\operatorname{Cov}(X_i,X_j)$.
\end{enumerate}
\hlt{Structural Commentary}: Expectation is linear; variance and covariance are quadratic and bilinear consequences.\\
\hlt{Key Insight}: Expand products, then use linearity.
\end{proof}
\end{theorem}

\begin{example} \label{ex:indicator-sums-binomial-random-walk-game} \hlt{Indicator Sums as a Method}
\begin{enumerate}[label=(\roman*)]
\setlength{\itemsep}{0pt}
\item If $X\sim\operatorname{Binomial}(n,p)$, write $X=\sum\limits_{i=1}^{n}I_i$ with independent $I_i\sim\operatorname{Bernoulli}(p)$.\\
Then $\mathbb{E}X=np$ and $\operatorname{Var}(X)=np(1-p)$.
\item Let $\xi_i\in\{-1,1\}$ with $\mathbb{P}(\xi_i=-1)=p$ and $\mathbb{P}(\xi_i=1)=1-p$, independently, and $X_n=\sum\limits_{i=1}^{n}\xi_i$.\\
Then $\mathbb{E}X_n=n(1-2p)$ and $\operatorname{Var}(X_n)=4np(1-p)$.
\item Let $F_n=c\prod\limits_{i=1}^{n}M_i$ with independent $M_i\in\{2,\frac{1}{2}\}$ and $\mathbb{P}(M_i=2)=\mathbb{P}(M_i=\frac{1}{2})=\frac{1}{2}$.\\
Then $\mathbb{E}F_n=c\prod\limits_{i=1}^{n}\mathbb{E}M_i=c(\frac{5}{4})^n$.
\end{enumerate}
\end{example}

\begin{proposition} \label{prop:variance-zero-iff-constant} \hlt{Zero Variance Means Almost Surely Constant}\\
For $X\in L^2$ and $m\defeq\mathbb{E}X$, $\operatorname{Var}(X)=0$ iff $X=m$ a.s.
\begin{proof}
\begin{enumerate}[label=\roman*.]
\setlength{\itemsep}{0pt}
\item[$(\Rightarrow)$] If $\operatorname{Var}(X)=0$, then $\mathbb{E}[(X-m)^2]=0$. Since $(X-m)^2\geq0$, $\mathbb{P}((X-m)^2>0)=0$, hence $X=m$ a.s.
\item[$(\Leftarrow)$] If $X=m$ a.s., then $(X-m)^2=0$ a.s., so $\operatorname{Var}(X)=\mathbb{E}[(X-m)^2]=0$.
\end{enumerate}
\hlt{Structural Commentary}: Variance is the mean squared distance from the mean.\\
\hlt{Key Insight}: Zero nonnegative expectation forces zero almost surely.
\end{proof}
\end{proposition}

\begin{proposition} \label{prop:tail-integral-formula-nonnegative} \hlt{Tail Integral Formula}\\
If $X\geq0$ has density $f_X$, distribution function $F_X$, and $\mathbb{E}X<\infty$,\\
then $\mathbb{E}X=\int_0^\infty\mathbb{P}(X>x)\,dx=\int_0^\infty(1-F_X(x))\,dx$.
\begin{proof}
For $b>0$, integration by parts gives $\int_0^bxf_X(x)\,dx=-b(1-F_X(b))+\int_0^b(1-F_X(x))\,dx$.\\
Since $\mathbb{E}X<\infty$, $b(1-F_X(b))\leq\int_b^\infty xf_X(x)\,dx\to0$.\\
Letting $b\to\infty$ yields $\mathbb{E}X=\int_0^\infty(1-F_X(x))\,dx=\int_0^\infty\mathbb{P}(X>x)\,dx$.\\
\hlt{Structural Commentary}: A nonnegative variable is integrated by summing its survival probabilities.\\
\hlt{Key Insight}: Expectation equals area under the tail.
\end{proof}
\end{proposition}

\begin{example} \label{ex:maximum-uniform-expectation} \hlt{Maximum of Uniform Samples}\\
Let $X_1,\ldots,X_n$ be independent $\operatorname{Uniform}(0,1)$ variables and $Y_n=\max_iX_i$.\\
Then $F_{Y_n}(y)=\mathbb{P}(X_1\leq y,\ldots,X_n\leq y)=y^n$ for $0\leq y\leq1$.\\
So $\mathbb{E}Y_n=\int_0^1(1-F_{Y_n}(y))\,dy=\int_0^1(1-y^n)\,dy=\frac{n}{n+1}$.
\end{example}

\begin{theorem} \label{thm:sample-mean-and-variance-moments} \hlt{Mean and Variance of the Basic Sample Statistics}\\
Let $X_1,\ldots,X_n$ be iid with $\mathbb{E}X_i=\mu$ and $\operatorname{Var}(X_i)=\sigma^2<\infty$.
\begin{enumerate}[label=(\roman*)]
\setlength{\itemsep}{0pt}
\item \hlt{Sample-mean expectation}: $\mathbb{E}\overline{X}_n=\mu$.
\item \hlt{Sample-mean variance}: $\operatorname{Var}(\overline{X}_n)=\frac{\sigma^2}{n}$.
\item \hlt{Unbiased sample variance}: $\mathbb{E}S_n^2=\sigma^2$.
\end{enumerate}
\begin{proof}
\begin{enumerate}[label=(\roman*)]
\setlength{\itemsep}{0pt}
\item $\mathbb{E}\overline{X}_n=n^{-1}\sum\limits_{i=1}^{n}\mathbb{E}X_i=n^{-1}n\mu=\mu$.
\item Independence gives $\operatorname{Var}(\overline{X}_n)=n^{-2}\sum\limits_{i=1}^{n}\operatorname{Var}(X_i)=n^{-2}n\sigma^2=\frac{\sigma^2}{n}$.
\item Since $\sum\limits_{i=1}^{n}(X_i-\overline{X}_n)^2=\sum\limits_{i=1}^{n}(X_i-\mu)^2-n(\overline{X}_n-\mu)^2$. Thus $\mathbb{E}[(n-1)S_n^2]=n\sigma^2-n(\frac{\sigma^2}{n})=(n-1)\sigma^2$.
\end{enumerate}
\hlt{Structural Commentary}: Averaging preserves mean and scales variance.\\
\hlt{Key Insight}: Center at $\mu$, then subtract the fluctuation of $\overline{X}_n$.
\end{proof}
\end{theorem}

\begin{definition} \label{def:random-vector-mean-covariance-matrix} \hlt{Mean Vector and Covariance Matrix}\\
For a random vector $X=(X_1,\ldots,X_k)^t$:
\begin{enumerate}[label=(\roman*)]
\setlength{\itemsep}{0pt}
\item \hlt{Mean vector}: $\mathbb{E}X\defeq(\mathbb{E}X_1,\ldots,\mathbb{E}X_k)^t$.
\item \hlt{Covariance matrix}: $\operatorname{Var}(X)\defeq[\operatorname{Cov}(X_i,X_j)]_{i,j=1}^{k}$.
\end{enumerate}
\end{definition}

\begin{proposition} \label{prop:linear-transform-random-vector-moments} \hlt{Linear Transformations of Random Vectors}\\
Let $X\in\R^k$ satisfy $\mathbb{E}X=\mu$ and $\operatorname{Var}(X)=\Sigma$.
\begin{enumerate}[label=\roman*.]
\setlength{\itemsep}{0pt}
\item \hlt{Linear form}: for $a\in\R^k$, $\mathbb{E}(a^tX)=a^t\mu$ and $\operatorname{Var}(a^tX)=a^t\Sigma a$.
\item \hlt{Matrix transform}: for $A\in\R^{m\times k}$, $\mathbb{E}(AX)=A\mu$ and $\operatorname{Var}(AX)=A\Sigma A^t$.
\end{enumerate}
\begin{proof}
\begin{enumerate}[label=\roman*.]
\setlength{\itemsep}{0pt}
\item $\mathbb{E}(a^tX)=a^t\mathbb{E}X=a^t\mu$, and $\operatorname{Var}(a^tX)=\mathbb{E}[(a^t(X-\mu))^2]=a^t\mathbb{E}[(X-\mu)(X-\mu)^t]a=a^t\Sigma a$.
\item $\mathbb{E}(AX)=A\mathbb{E}X=A\mu$, and $\operatorname{Var}(AX)=\mathbb{E}[(AX-A\mu)(AX-A\mu)^t]=A\mathbb{E}[(X-\mu)(X-\mu)^t]A^t=A\Sigma A^t$.
\end{enumerate}
\hlt{Structural Commentary}: Linear maps push means forward and conjugate covariance matrices.\\
\hlt{Key Insight}: Center first; then move constants outside expectations.
\end{proof}
\end{proposition}

\begin{example} \label{ex:multinomial-covariance-matrix} \hlt{Multinomial Covariance Structure}\\
If $X\sim\operatorname{Multinomial}(n,p)$ with $p=(p_1,\ldots,p_k)$, then $\mathbb{E}X=np$ and $\operatorname{Var}(X_i)=np_i(1-p_i)$.\\
For $i\neq j$, $X_i+X_j\sim\operatorname{Binomial}(n,p_i+p_j)$, so $n(p_i+p_j)(1-p_i-p_j)=np_i(1-p_i)+np_j(1-p_j)+2\operatorname{Cov}(X_i,X_j)$.\\
Hence $\operatorname{Cov}(X_i,X_j)=-np_ip_j$.
\end{example}

\begin{theorem} \label{thm:markov-and-chebyshev-inequalities} \hlt{Markov and Chebyshev Inequalities}
\begin{enumerate}[label=\roman*.]
\setlength{\itemsep}{0pt}
\item \hlt{Markov}: if $X\geq0$ and $\mathbb{E}X<\infty$, then for $t>0$, $\mathbb{P}(X\geq t)\leq\frac{\mathbb{E}X}{t}$.
\item \hlt{Chebyshev}: if $\mu\defeq\mathbb{E}X$ and $\sigma^2\defeq\operatorname{Var}(X)<\infty$, then for $t>0$, $\mathbb{P}(\abs{X-\mu}\geq t)\leq\frac{\sigma^2}{t^2}$.
\end{enumerate}
\begin{proof}
\begin{enumerate}[label=\roman*.]
\setlength{\itemsep}{0pt}
\item Since $X\geq t\mathbf{1}_{\{X\geq t\}}$, taking expectations gives $\mathbb{E}X\geq t\mathbb{P}(X\geq t)$.
\item Apply (i) to $Y\defeq(X-\mu)^2\geq0$ with threshold $t^2$, then $\mathbb{P}(\abs{X-\mu}\geq t)=\mathbb{P}(Y\geq t^2)\leq\frac{\mathbb{E}Y}{t^2}=\frac{\sigma^2}{t^2}$.
\end{enumerate}
\hlt{Structural Commentary}: Tail bounds convert moment control into probability control.\\
\hlt{Key Insight}: Dominate the tail indicator by a scaled nonnegative random variable.
\end{proof}
\end{theorem}

\begin{example} \label{ex:chebyshev-error-rate-bound} \hlt{A Distribution-Free Error-Rate Bound}\\
Let $X_1,\ldots,X_n$ be independent $\operatorname{Bernoulli}(p)$ error indicators and $\overline{X}_n\defeq n^{-1}\sum\limits_{i=1}^{n}X_i$. Since $\mathbb{E}\overline{X}_n=p$ and $\operatorname{Var}(\overline{X}_n)=\frac{p(1-p)}{n}$, Chebyshev gives, for $\varepsilon>0$, $\mathbb{P}(\abs{\overline{X}_n-p}>\varepsilon)\leq\frac{p(1-p)}{n\varepsilon^2}\leq\frac{1}{4n\varepsilon^2}$.
\end{example}

\begin{theorem} \label{thm:hoeffding-bounded-zero-mean} \hlt{Hoeffding's Bounded-Sum Bound}\\
Let $Y_1,\ldots,Y_n$ be independent, $\mathbb{E}Y_i=0$, $a_i\leq Y_i\leq b_i$ a.s., and $V\defeq\sum\limits_{i=1}^{n}(b_i-a_i)^2$. For $S_n\defeq\sum\limits_{i=1}^{n}Y_i$ and $t,\varepsilon>0$, $\mathbb{P}(S_n\geq\varepsilon)\leq\exp(-t\varepsilon+\frac{t^2V}{8})$; if $V>0$, then $\mathbb{P}(S_n\geq\varepsilon)\leq\exp(-\frac{2\varepsilon^2}{V})$.
\begin{proof}
For $t>0$, Markov gives $\mathbb{P}(S_n\geq\varepsilon)=\mathbb{P}(e^{tS_n}\geq e^{t\varepsilon})\leq e^{-t\varepsilon}\mathbb{E}e^{tS_n}$.\\
Independence gives $\mathbb{E}e^{tS_n}=\prod\limits_{i=1}^{n}\mathbb{E}e^{tY_i}$, and Hoeffding's lemma gives $\mathbb{E}e^{tY_i}\leq\exp(\frac{t^2(b_i-a_i)^2}{8})$.\\
Hence $\mathbb{P}(S_n\geq\varepsilon)\leq\exp(-t\varepsilon+\frac{t^2V}{8})$. If $V>0$, minimize the exponent at $t=\frac{4\varepsilon}{V}$ to obtain $\exp(-\frac{2\varepsilon^2}{V})$.\\
\hlt{Structural Commentary}: Hoeffding turns bounded independent increments into exponential tail decay.\\
\hlt{Key Insight}: Exponentiate the sum, factor the MGF, then optimize $t$.
\end{proof}
\end{theorem}

\begin{corollary} \label{cor:hoeffding-bernoulli-average} \hlt{Hoeffding for Bernoulli Averages}\\
If $X_1,\ldots,X_n$ independent $\operatorname{Bernoulli}(p)$ variables, $\overline{X}_n\defeq n^{-1}\sum\limits_{i=1}^{n}X_i$, then $\mathbb{P}(\abs{\overline{X}_n-p}>\varepsilon)\leq2e^{-2n\varepsilon^2}$.
\begin{proof}
Set $Y_i^+\defeq\frac{X_i-p}{n}$.\\
Then $\mathbb{E}Y_i^+=0$, $Y_i^+\in[-\frac{p}{n},\frac{1-p}{n}]$, and $V=\sum\limits_i(\frac{1}{n})^2=\frac{1}{n}$, by Theorem \ref{thm:hoeffding-bounded-zero-mean}, $\mathbb{P}(\overline{X}_n-p\geq\varepsilon)\leq e^{-2n\varepsilon^2}$.\\
Apply same argument to $Y_i^-\defeq\frac{p-X_i}{n}$, $\mathbb{P}(p-\overline{X}_n\geq\varepsilon)\leq e^{-2n\varepsilon^2}$; the union bound yields the claim.\\
\hlt{Structural Commentary}: Bernoulli averages are sums of bounded centered increments.\\
\hlt{Key Insight}: Control both tails separately, then add them.
\end{proof}
\end{corollary}

\begin{theorem} \label{thm:cauchy-schwarz-and-jensen-expectations} \hlt{Cauchy--Schwarz and Jensen}\\
If $X,Y\in L^2$, then $\mathbb{E}\abs{XY}\leq\sqrt{\mathbb{E}(X^2)\mathbb{E}(Y^2)}$.\\
If $g:I\to\mathbb{R}$ convex, $X\in I$ a.s., $X,g(X)\in L^1$, then $\mathbb{E}g(X)\geq g(\mathbb{E}X)$; for concave $g$, inequality reverses.
\begin{proof}
For Cauchy--Schwarz, apply the $L^2$ inner-product inequality to $\abs{X}$ and $\abs{Y}$.\\
For Jensen, set $m\defeq\mathbb{E}X$ and take an affine support $\ell(x)\defeq g(m)+s(x-m)$ with $\ell\leq g$ on $I$.\\
Then $\mathbb{E}g(X)\geq\mathbb{E}\ell(X)=g(m)+s(\mathbb{E}X-m)=g(\mathbb{E}X)$. For concave $g$, apply the convex case to $-g$.\\
\hlt{Structural Commentary}: Cauchy--Schwarz controls products; Jensen moves convexity through expectation.\\
\hlt{Key Insight}: Use Hilbert-space geometry for products and supporting lines for convex functions.
\end{proof}
\end{theorem}

\begin{definition} \label{def:moment-generating-function} \hlt{Moment Generating Function}\\
The \hlt{moment generating function} of $X$ is $\psi_X(t)\defeq\mathbb{E}(e^{tX})$ on $D_X\defeq\{t\in\mathbb{R}:\mathbb{E}e^{tX}<\infty\}$.\\
If $(-\delta,\delta)\subseteq D_X$ for some $\delta>0$, then $\psi_X^{(k)}(0)=\mathbb{E}(X^k)$ for all $k\geq1$.
\end{definition}

\begin{proposition} \label{prop:mgf-basic-properties} \hlt{Basic MGF Properties}\\
Assume the MGFs below are finite on a common open interval about $0$.
\begin{enumerate}[label=(\roman*)]
\setlength{\itemsep}{0pt}
\item \hlt{Affine transform}: if $a,b\in\mathbb{R}$ and $Y=aX+b$, then $\psi_Y(t)=e^{bt}\psi_X(at)$.
\item \hlt{Independent sum}: if $X_1,\ldots,X_n$ are independent and $Y\defeq\sum\limits_{i=1}^{n}X_i$, then $\psi_Y(t)=\prod\limits_{i=1}^{n}\psi_{X_i}(t)$.
\item \hlt{Uniqueness}: if $\psi_X(t)=\psi_Y(t)$ for all $t$ in an open interval containing $0$, then $\mathcal{L}(X)=\mathcal{L}(Y)$.
\end{enumerate}
\begin{proof}
\begin{enumerate}[label=(\roman*)]
\setlength{\itemsep}{0pt}
\item $\psi_Y(t)=\mathbb{E}e^{t(aX+b)}=e^{bt}\mathbb{E}e^{(at)X}=e^{bt}\psi_X(at)$.
\item For admissible $t$, $e^{tX_1},\ldots,e^{tX_n}$ are independent and integrable.\\
Hence $\psi_Y(t)=\mathbb{E}e^{t\sum\limits_{i=1}^{n}X_i}=\mathbb{E}\prod\limits_{i=1}^{n}e^{tX_i}=\prod\limits_{i=1}^{n}\mathbb{E}e^{tX_i}=\prod\limits_{i=1}^{n}\psi_{X_i}(t)$.
\item Let $\mu\defeq\mathcal{L}(X)$ and $\nu\defeq\mathcal{L}(Y)$. If $\psi_X=\psi_Y$ on $(-\delta,\delta)$, then $\int_{\mathbb{R}}e^{tx}\,d\mu(x)=\int_{\mathbb{R}}e^{tx}\,d\nu(x)$ for all $\abs{t}<\delta$.\\
Uniqueness of the two-sided Laplace transform gives $\mu=\nu$.
\end{enumerate}
\hlt{Structural Commentary}: MGFs turn affine maps and independent sums into algebraic operations.\\
\hlt{Key Insight}: Exponentials convert sums into products, and the MGF near $0$ determines the law.
\end{proof}
\end{proposition}

\begin{example} \label{ex:mgf-computations} \hlt{MGF Computations}
\begin{enumerate}[label=(\roman*)]
\setlength{\itemsep}{0pt}
\item If $X\sim\operatorname{Exp}(1)$, then $\psi_X(t)=\mathbb{E}e^{tX}=\int_{0}^{\infty}e^{(t-1)x}\,dx=(1-t)^{-1}$ for $t<1$.\\
Hence $\mathbb{E}X=\psi_X'(0)=1$ and $\operatorname{Var}(X)=\psi_X''(0)-(\psi_X'(0))^2=1$.
\item If $X\sim\operatorname{Binomial}(n,p)$, then $X=\sum\limits_{i=1}^{n}B_i$ with independent $B_i\sim\operatorname{Bernoulli}(p)$.\\
Hence $\psi_X(t)=\prod\limits_{i=1}^{n}\psi_{B_i}(t)=(pe^t+1-p)^n$.
\end{enumerate}
\end{example}

\begin{method} \label{meth:common-moments-and-mgfs-table} \hlt{Common Means, Variances, and MGFs}\\
MGFs are stated on their natural neighborhood of $0$; values at $t=0$ are by continuity.\\
The geometric and negative-binomial rows use the total-trials convention.
\[
\begin{array}{c|c|c|c}
\text{Distribution} & \mathbb{E}X & \operatorname{Var}(X) & \psi_X(t)\\ \hline
\text{point mass at }a & a & 0 & e^{at}\\
\operatorname{Bernoulli}(p) & p & p(1-p) & 1-p+pe^t\\
\operatorname{Binomial}(n,p) & np & np(1-p) & (1-p+pe^t)^n\\
\operatorname{Geometric}(p) & \dfrac{1}{p} & \dfrac{1-p}{p^2} & \dfrac{pe^t}{1-(1-p)e^t}\\
\operatorname{NegBin}(r,p) & \dfrac{r}{p} & \dfrac{r(1-p)}{p^2} & \left(\dfrac{pe^t}{1-(1-p)e^t}\right)^r\\
\operatorname{Poisson}(\lambda) & \lambda & \lambda & e^{\lambda(e^t-1)}\\
\operatorname{Uniform}(a,b) & \dfrac{a+b}{2} & \dfrac{(b-a)^2}{12} & \dfrac{e^{bt}-e^{at}}{(b-a)t}\\
N(\mu,\sigma^2) & \mu & \sigma^2 & \exp(\mu t+\dfrac{\sigma^2t^2}{2})\\
\operatorname{Exp}(\lambda) & \dfrac{1}{\lambda} & \dfrac{1}{\lambda^2} & \dfrac{\lambda}{\lambda-t}\\
\operatorname{Gamma}(\alpha,\lambda) & \dfrac{\alpha}{\lambda} & \dfrac{\alpha}{\lambda^2} & \left(\dfrac{\lambda}{\lambda-t}\right)^\alpha\\
\operatorname{Beta}(\alpha,\beta) & \dfrac{\alpha}{\alpha+\beta} & \dfrac{\alpha\beta}{(\alpha+\beta)^2(\alpha+\beta+1)} & \text{not useful}\\
\chi_p^2 & p & 2p & (1-2t)^{-\frac{p}{2}}
\end{array}
\]
\end{method}
